\section{\texorpdfstring{光滑经验风险与 \(L_2\) 正则化}{光滑经验风险与 L2 正则化}}
\label{sec:09-Convex-Optimization/07-Applications/01}

特征表里有两列几乎一样：一列叫 \texttt{spend\_7d}，一列叫 \texttt{spend\_7d\_fixed}，是数据修复时补录的同一个量，相关系数 \(0.998\)。线性回归跑出来，这两列的系数是 \(312\) 和 \(-309\)。换一个随机划分重跑，变成 \(-147\) 和 \(150\)。训练误差两次几乎相同，报告里的``该特征对留存的影响''却换了正负号。

问题不在求解器。目标函数在``两个系数同增同减、和保持不变''这个方向上几乎是平的，沿着这条方向走出很远，损失也只掉了小数点后第五位。求解器每次停在这条平谷上的哪一点，取决于初值、舍入和数据的微小抖动。它交出的是一个最优解，不是那个最优解。

在目标上加一项 \(\frac{\lambda}{2}\norm{w}_2^2\)，两次重跑的系数变成 \(1.61\) 和 \(1.58\)。这一节要讲清楚这一项做了什么，以及为什么``防过拟合''这个说法只说中了它一半的作用。

\subsection{学习先被写成一个优化问题}

监督学习的标准写法是经验风险最小化：手上有 \(n\) 个样本 \((x_i,y_i)\)，用一族由参数 \(w\) 索引的预测函数去拟合，把``拟合得好不好''量化成损失 \(\ell\)，再对样本取平均：

\[
\min_{w\in\R^d}\ \frac1n\sum_{i=1}^n\ell\bigl(y_i,\,w^{\trans}x_i\bigr)+\lambda R(w).
\]

线性预测 \(w^{\trans}x_i\) 是骨架，\(\ell\) 决定模型怎样看待一次预测失误，\(R\) 决定在同样拟合水平下偏好什么样的参数。这一步就是\hyperref[sec:09-Convex-Optimization/01-Optimization-Foundations/01]{``把现实选择翻译成优化问题''}里说的翻译动作：真正关心的量是未来数据上的期望损失，写下来的却是手头这 \(n\) 个样本上的平均。两者的差距后面会算账，眼下先看这个已经写定的目标长什么样。

若 \(\ell\) 对第二个变量凸，则它与仿射映射 \(w\mapsto w^{\trans}x_i\) 复合后关于 \(w\) 仍凸，求和取平均保持凸，加上凸的 \(R\) 还是凸。平方损失、logistic 损失、hinge 损失都满足前提，所以整族问题一次性落进凸优化。

没有正则项时，目标的曲率完全由数据提供。平方损失的 Hessian 是 \(\frac1nX^{\trans}X\)，其中 \(X\) 是把 \(x_i^{\trans}\) 按行堆起来的设计矩阵。开头那两列共线的特征，让这个矩阵有一个近乎为零的特征值，对应的特征向量正是那条平谷的方向。若 \(d>n\)，零特征值至少有 \(d-n\) 个，平谷变成一整个仿射子空间。

\subsection{一项 \texorpdfstring{\(\lambda I\)}{lambda I} 把平谷抬成盆地}

Ridge 回归就是在平方损失上加 \(L_2\) 惩罚：

\[
\min_{w}\ \frac{1}{2n}\norm{Xw-y}_2^2+\frac{\lambda}{2}\norm{w}_2^2 .
\]

系数 \(1/(2n)\) 让损失是每样本平均，\(\lambda/2\) 让惩罚项的梯度恰好是 \(\lambda w\)。Hessian 变成

\[
H_\lambda=\frac1nX^{\trans}X+\lambda I .
\]

原来的特征值 \(\mu_1\ge\cdots\ge\mu_d\ge0\) 整体平移成 \(\mu_j+\lambda\)。曲率大的方向几乎不受影响，因为 \(\lambda\) 相对 \(\mu_1\) 很小；曲率接近零的方向被抬到 \(\lambda\)。

\begin{center}
\begin{tikzpicture}[font=\small,>=stealth]
  \node[font=\footnotesize,black!70] at (0,2.15) {\(\lambda=0\)：一条平谷};
  \begin{scope}[rotate=32]
    \draw[black!70] (0,0) ellipse (2.05 and 0.20);
    \draw[black!55] (0,0) ellipse (1.45 and 0.14);
    \draw[black!40] (0,0) ellipse (0.80 and 0.08);
    \draw[dashed,black!45] (-2.35,0) -- (2.35,0);
    \fill[black!75] (-1.75,0) circle (1.6pt);
    \fill[black!75] (-0.55,0) circle (1.6pt);
    \fill[black!75] (1.30,0) circle (1.6pt);
  \end{scope}
  \node[font=\footnotesize,black!60] at (-1.35,-1.75) {三次重跑，三个答案};
  \begin{scope}[shift={(6.2,0)}]
    \node[font=\footnotesize,black!70] at (0,2.15) {\(\lambda>0\)：一个盆地};
    \begin{scope}[rotate=32]
      \draw[black!70] (0,0) ellipse (1.45 and 0.86);
      \draw[black!55] (0,0) ellipse (1.00 and 0.59);
      \draw[black!40] (0,0) ellipse (0.55 and 0.33);
      \fill[black!85] (0,0) circle (2.0pt);
    \end{scope}
    \node[font=\footnotesize,black!60] at (0,-1.75) {三次重跑，同一个答案};
  \end{scope}
\end{tikzpicture}

\emph{同一份数据的等值线：左边沿细长方向几乎没有曲率，谷底是一条线；\(\lambda I\) 给每个方向补上同样大小的曲率，等值线变圆，谷底收缩成一点。}
\end{center}

图上左右两组等值线来自同一个 \(X\)。左边最长的那根轴，半长与 \(1/\sqrt{\mu_d}\) 成正比，\(\mu_d\) 趋于零它就趋于无穷，谷底退化成一条直线，直线上每一点的训练损失都是最小值。右边的半长变成 \(1/\sqrt{\mu_d+\lambda}\)，有限了，于是谷底只剩一个点。开头那两次重跑给出 \(312\) 与 \(-147\)，就是求解器停在左图那条虚线上的两个不同位置。

\subsection{数值这笔账：唯一、有速率、可复算}

\(H_\lambda\) 的最小特征值至少是 \(\lambda\)，这正是 \(\lambda\)-强凸的定义，见\hyperref[sec:09-Convex-Optimization/03-Convex-Functions/03]{``曲率区分凸、严格凸与强凸''}。那一篇里的三条结论此刻同时生效。

最优解唯一，且落在有限处。强凸函数的最优解至多一个，而 \(\norm{w}\to\infty\) 时目标趋于无穷，强制性保证它至少有一个。于是``系数是多少''这个问题有了单一答案，重跑得到同一个数，代码审查和结果复现才有意义。

收敛有速率。梯度下降在 \(\lambda\)-强凸、\(L\)-光滑目标上以 \(1-\lambda/L\) 的比例线性收敛，迭代次数正比于条件数 \(\kappa=L/\lambda\)，见\hyperref[sec:09-Convex-Optimization/06-Optimization-Algorithms/01]{``梯度下降真正难在步长''}。这里 \(L=\mu_1+\lambda\)，所以 \(\kappa=(\mu_1+\lambda)/(\mu_d+\lambda)\)：不加正则时 \(\mu_d\approx0\) 让 \(\kappa\) 爆炸，加了以后 \(\kappa\le(\mu_1+\lambda)/\lambda\)，与数据的病态程度脱钩。共轭梯度的迭代次数依赖 \(\sqrt{\kappa}\)，同样受益。

解对数据扰动不敏感。一阶条件给出正规方程 \(H_\lambda w=\frac1nX^{\trans}y\)。把响应 \(y\) 换成 \(y+\delta\)，解的变动满足

\[
\begin{aligned}
H_\lambda\,\Delta w&=\tfrac1nX^{\trans}\delta,\\
\norm{\Delta w}_2&\le\norm{H_\lambda^{-1}}_2\cdot\tfrac1n\norm{X^{\trans}\delta}_2
\le\frac{\norm{X^{\trans}\delta}_2}{n\lambda}.
\end{aligned}
\]

骨架三步：正规方程对 \(y\) 是线性的，所以扰动可以整体移到右端；\(H_\lambda\) 的谱范数的逆等于最小特征值的倒数；最小特征值不小于 \(\lambda\)。条件对账：第一步用到平方损失（换成 logistic 损失只有局部的线性化版本），第二步用到 \(H_\lambda\) 对称，第三步用到 \(X^{\trans}X\succeq0\) 与 \(\lambda>0\)。全程没有假设 \(X\) 列满秩——这正是重点，\(\lambda\) 把满秩这个前提买了下来。灵敏度的一般语言见\hyperref[sec:08-Numerical-Analysis/01-Floating-Point-and-Error/03]{``条件数衡量问题本身的敏感性''}。

实现上还有两句话。正规方程要解而不要求逆，Cholesky（\(H_\lambda\) 对称正定）、QR（直接作用在 \(X\) 上更稳）、SVD（一次分解扫遍所有 \(\lambda\)）按规模挑一个，理由见\hyperref[sec:08-Numerical-Analysis/07-Numerical-Methods-for-ML/03]{``需要的是解，而不是逆矩阵''}。另有一个等价的增广写法

\[
\min_w\ \frac{1}{2n}\norm{\begin{bmatrix}X\\ \sqrt{n\lambda}\,I\end{bmatrix}w-\begin{bmatrix}y\\ 0\end{bmatrix}}_2^2,
\]

它把正则化重新表述成一个普通最小二乘：往数据里塞 \(d\) 条虚拟样本，每条都说``这个系数应该是零''，权重 \(\sqrt{n\lambda}\)。任何现成的最小二乘求解器可以直接复用。

\subsection{统计这笔账：用一点偏差换一大块方差}

上面几条全部是关于计算的：唯一性、迭代次数、舍入敏感度。它们成立与否，和数据是怎么生成的、样本够不够、模型对不对，一概无关。统计上的收益要另外算。

Ridge 解可以写成 \(\hat w_\lambda=(X^{\trans}X+n\lambda I)^{-1}X^{\trans}y\)。在 \(X\) 的奇异值分解下逐个奇异方向看，第 \(j\) 个方向的最小二乘估计被乘上收缩因子 \(s_j^2/(s_j^2+n\lambda)\)。奇异值大的方向因子接近 \(1\)，基本不动；奇异值小的方向因子接近 \(0\)，被压向零。而恰恰是小奇异值方向上，估计的方差正比于 \(1/s_j^2\)，噪声被放得最大。\(\lambda\) 做的事是：在数据支撑最弱的方向上，主动放弃一部分信号，换掉一大部分噪声。估计因此有了偏差，方差因此下降，总的均方误差在某个 \(\lambda^\star>0\) 处取到最小——这条 U 形曲线是\hyperref[sec:11-Mathematical-Statistics/06-Resampling-and-Model-Selection/04]{``偏差---方差、正则化与模型选择''}的主题。

两笔账的方向一致，但不能合并。数值那一侧，\(\lambda\) 越大越好：条件数越小，收敛越快，解越稳。统计那一侧，\(\lambda\) 太大会把有用的信号一起压掉，欠拟合。所以 \(\lambda\) 的取值必须由统计那一侧决定——用独立验证集或交叉验证选，看验证损失的 U 形底点——数值上的好处只是搭了个便车。反过来说，如果有人为了让求解器收敛快而调大 \(\lambda\)，他是在用模型精度补贴计算预算，这个交换应该被明说。

还有两处工程细节值得单列。惩罚 \(\norm{w}_2^2\) 对特征尺度不中立：同一个变量以米为单位和以毫米为单位，系数差三个数量级，受到的惩罚差六个数量级。所以标准化必须在正则化之前完成，而且标准化的均值和方差只能从训练折里算，否则就是数据泄漏。截距通常不该惩罚，它不乘任何特征，压它等于强迫拟合曲面穿过原点附近。

\subsection{Logistic 回归里正则还兼职保证解的存在}

换成分类问题，标签取 \(y_i\in\set{-1,1}\)，记间隔 \(t_i=y_i(w^{\trans}x_i+b)\)，logistic 损失是 \(\ell(t)=\log(1+e^{-t})\)。它的二阶导数 \(\ell''(t)=\sigma(t)\bigl(1-\sigma(t)\bigr)\) 处处为正，与仿射映射复合后对 \(w\) 凸。整个训练目标的 Hessian 是

\[
\frac1nX^{\trans}WX+\lambda I,\qquad W=\diag\bigl(p_1(1-p_1),\dots,p_n(1-p_n)\bigr),
\]

其中 \(p_i\) 是模型给样本 \(i\) 的预测概率。权重 \(p_i(1-p_i)\) 在 \(p_i=1/2\) 处取最大值 \(1/4\)，在 \(p_i\) 接近 \(0\) 或 \(1\) 时趋于零。这句话的含义值得停一下：曲率几乎全部来自决策边界附近那些模型拿不准的样本。被远远甩在边界外侧的多数类样本，在损失值和梯度里占大头，在 Hessian 里几乎不出现。类别严重不平衡时，有效曲率来自一小撮样本，条件数比样本量看上去暗示的要糟得多。Newton 法与迭代重加权最小二乘每一步解的就是这个加权系统，权重随预测概率自适应地重新分配。

不加正则时，logistic 回归有个只在凸性框架里才看得清的毛病。若训练数据线性可分，那么把任意一个分对全部样本的 \(w\) 乘以常数 \(c\to\infty\)，所有间隔同步放大，损失单调趋于下确界 \(0\)，却在任何有限的 \(w\) 处都取不到。凸性保证局部最优即全局最优，从不保证最优解存在——这一点在\hyperref[sec:09-Convex-Optimization/01-Optimization-Foundations/02]{``局部更好为何还不够''}里已经交代过。加上 \(\lambda\norm{w}_2^2\) 后强制性恢复，最优解在有限处唯一存在。工程上的症状很好认：不加正则又恰好可分时，训练损失一直缓慢下降，参数范数一路增长，梯度范数永远不到停机阈值，日志上看不出任何异常，作业就是不结束。

数值稳定还需要一句提醒。\(\log(1+e^{z})\) 在 \(z=100\) 时直接溢出，应改写成 \(\max(0,z)+\log(1+e^{-\abs{z}})\)；同样不要先算出概率再取对数，\(\sigma(z)\) 在 \(z\) 很负时下溢成零。这类改写的通则见\hyperref[sec:08-Numerical-Analysis/07-Numerical-Methods-for-ML/01]{``在对数域中稳定计算概率''}。

\subsection{惩罚形式与约束形式说的是同一件事}

\(L_2\) 正则还有一种写法，把``系数别太大''直接放进可行域：

\[
\begin{aligned}
\text{minimize}\quad & \ell(w)\\
\text{subject to}\quad & \norm{w}_2^2\le t .
\end{aligned}
\]

在凸性与 Slater 条件下，这个问题与惩罚形式 \(\min_w \ell(w)+\lambda\norm{w}_2^2\) 通过 Lagrange 乘子一一对应：每个半径 \(t\) 配一个乘子 \(\lambda(t)\ge0\)，它等于最优值对 \(t\) 的导数的相反数，也就是``把预算放宽一个单位能省下多少损失''，即\hyperref[sec:09-Convex-Optimization/05-Duality/04]{``对偶变量是约束的影子价格''}里的影子价格。这个对应是隐式的，\(\lambda\) 不是 \(1/t\) 之类的现成公式，它依赖数据。约束不活跃时——无正则最优解本来就落在球内——对应的 \(\lambda=0\)，约束等于没写。

哪种写法方便取决于场景。参数范数有物理预算时（``总能量不超过某值''）约束形式更自然；要用交叉验证扫一条正则化路径时惩罚形式更方便，因为 \(\lambda\) 的网格可以取对数等距，而且很多求解器支持从上一个 \(\lambda\) 的解热启动。

\subsection{优化收敛不等于学习完成}

KKT 残差降到 \(10^{-9}\)、对偶间隙小于 \(0.1\%\)，说明写下来的那个目标已经被压到接近全局最优。它没说的是这个目标本身是否值得压。验证损失可能在训练损失还在下降时已经转头向上，概率输出可能系统性地过度自信，特征表里可能混进了预测时刻拿不到的字段。这些都不是优化误差，求解器再精确一分也不会改善半点。把优化误差、统计误差与分布漂移三者拆开归因，是\hyperref[sec:09-Convex-Optimization/08-Nonconvex-and-Stochastic-Optimization/06]{``优化误差、统计误差与泛化不能混为一谈''}的内容。

回到开头那两列共线的特征。\(L_2\) 正则让系数变成 \(1.61\) 和 \(1.58\)，稳定了，可它没有回答``哪一列该留下''这个问题——它把权重在两列之间平摊了。收缩因子 \(s_j^2/(s_j^2+n\lambda)\) 永远严格为正，所以 \(L_2\) 会把系数压小，几乎从不把任何一个压成恰好为零。当交付要求是``从三百个特征里挑出十个上线''时，光靠收缩不够，得换一种惩罚的形状。

那么把 \(\norm{w}_2^2\) 换成 \(\norm{w}_1\) 会发生什么？答案与其说在统计里，不如说在几何里：球换成了菱形，而菱形有尖角。下一篇\hyperref[sec:09-Convex-Optimization/07-Applications/02]{``\(L_1\) 正则为何产生稀疏''}从这个尖角开始。
